\documentclass[10pt,a4paper]{amsart}
\usepackage[margin=19mm]{geometry}
\usepackage{amsmath,amssymb,mathtools}
\usepackage[scr=rsfs,cal=euler]{mathalfa}
\usepackage{xfrac}
\usepackage[unicode,hidelinks,bookmarks=false]{hyperref}
\newcommand{\ie}{i.e.\ }
\newcommand{\cf}{cf.\ }
\newcommand{\resp}{respectively\ }
\newcommand{\Subsection}{\subsection}
\providecommand{\nobreakdash}{\nobreak\hbox{-}}
\setlength{\emergencystretch}{2em}
\multlinegap0pt
\usepackage{amssymb}
\usepackage{xcolor}
\usepackage{mathtools}
\usepackage[shortlabels]{enumitem}
\usepackage{tcolorbox}
\newtheorem{theorem}{Theorem}
\title{On the geometry of $G$-norm}
\author{Lakshmi Kanta Dey, Subhadip Pal}
\date{Source: arXiv:2603.20796v1 (CC BY 4.0)}
\begin{document}
\maketitle

\noindent\emph{Source and licence.} On the geometry of $G$-norm. Version arXiv:2603.20796v1 (CC BY 4.0), distributed by arXiv under the Creative Commons Attribution 4.0 International licence, http://creativecommons.org/licenses/by/4.0/, as recorded in the arXiv licence metadata for this exact version and shown on the abstract page https://arxiv.org/abs/2603.20796v1. Full licence notice: \texttt{licenses/arxiv-2603.20796-CC-BY-4.0.txt}.\par\smallskip

\noindent\textit{Editorial scope.} Counted unit: the article's theorem Th3.12 (source0.tex lines 715--727). The article's complete original proof of it is delivered with it (source0.tex lines 728--769, one \texttt{proof} environment of 42 lines). No mathematics is written, repaired or re-derived: the statement and the proof are the article's own lines, in the article's own notation, and no retained line is substituted or re-flowed.  No further source-local context is needed: every cross-reference of every retained line resolves inside the two delivered spans.  Nothing was omitted from any retained span, so the package carries no omission record.  No retained line contains a citation, so the wrapper carries no bibliography and no bibliography entry.  No retained line contains a figure environment, an image-inclusion command or drawing code, so no image is delivered and the package declares no asset dependency.  Editorial formatting only, in addition: an AMS \texttt{amsart} preamble that restates the article's own declarations and macros used by the retained lines, namely the environments \texttt{theorem} on separate counters, together with the standard packages those lines need.  The retained environments print in this document as Theorem 1 in order of appearance, whereas the article prints the same items with the numbering of its own full document.  The complete native source of the article as distributed in the arXiv e-print for this exact version is bundled unchanged as \texttt{sources/196878/source0.tex}.
	\begin{theorem}\label{Th3.12}
		Let $H$ be a Hilbert space and let 
		$G \in L(H)$ be a norm-attaining operator with $\|G\|=1$. Let $E = \overline{\operatorname{span}}(M_G).$ Then the following are equivalent:
		\begin{enumerate}
			\item[(i)] $\gamma := \sup\{\|Gx\| : x \in S_H \cap E^\perp\} < 1.$
			\item[(ii)] For every $\varepsilon>0$ there exists $\delta>0$ such that
			\[
			x\in S_H,\ \|Gx\|>1-\delta \Longrightarrow \operatorname{dist}(x,E)<\varepsilon.
			\]
			\item[(iii)] For every $T\in L(H)$,
			$\|T\|_G = \sup_{x\in M_G}\|Tx\|.$
		\end{enumerate}
	\end{theorem}

	\begin{proof}
		First, we establish two properties of $G$ on $E$. Since $\|G\|=1$, the operator $I-G^*G$ is positive. For $x\in M_G$ we have $\langle G^*Gx,x\rangle=\|Gx\|^2=1=\langle x, x\rangle.$ Consequently, $\langle (I-G^*G)x,x\rangle=0$. Positivity implies $(I-G^*G)x=0$, so $G^*Gx=x$. By continuity this extends to $E$, and therefore $G^*G=I$ on $E.$ Consequently,
		\begin{enumerate}
			\item[(a)] $G$ is an isometry on $E$ and
			\item[(b)] if $u\in E$ and $v\in E^\perp$, then $\langle Gu,Gv\rangle=0$.
		\end{enumerate}
		Since $E$ is a closed subspace of the Hilbert space $H$, there exists an orthogonal projection $P:H\to E$. For $x\in S_H$ we write $x = Px + (I-P)x,$ where $Px\in E$ and $(I-P)x\in E^\perp$. Let $u=Px$ and $v=(I-P)x$. Then by using the fact in (a) and (b) we have
		\begin{equation}\label{eq1}
			\|Gx\|^2=\|u\|^2+\|Gv\|^2 .
		\end{equation}
		
		
		\medskip
		\noindent
		$(i)\implies (ii)$. Assume $\|Gv\|\le \gamma\|v\|$ for $v\in E^\perp$ with $\gamma<1$. Let $\|x\|=1$ then $x=u+v$ for some $u\in E$ and $v\in E^\perp$. Therefore, $u\perp v$ implies $1=\|u\|^2+\|v\|^2$. It then follows from \eqref{eq1} that $\|Gx\|^2\le 1-(1-\gamma^2)\|v\|^2.$ If $\|Gx\|>1-\delta$, then $(1-\delta)^2<1-(1-\gamma^2)\|v\|^2,$ so $\|v\|^2<\frac{2\delta-\delta^2}{1-\gamma^2}.$ It is easy to see that $\operatorname{dist}(x, E)=\|v\|$, thus, by choosing $\delta$ sufficiently small gives $\operatorname{dist}(x,E)<\varepsilon$.
		
		
		
		\medskip
		
		\noindent
		$(ii)\implies (iii).$ Since $M_G \subseteq \{x\in S_H:\|Gx\|>1-\delta\}$, for every $\delta>0$, $\sup_{\|Gx\|>1-\delta}\|Tx\|\ge \sup_{x\in M_G}\|Tx\|.$ Consequently, $\|T\|_G \ge \sup_{x\in M_G}\|Tx\|$. If we fix $\varepsilon>0$, there exists $\delta>0$ such that
		\begin{equation}\label{eq2}
			\|Gx\|>1-\delta \implies \operatorname{dist}(x,E)<\varepsilon.
		\end{equation}
		Choose $z\in E$ with $\|x-z\|<\varepsilon$ and set $y=z/\|z\|\in M_G$. Then $\|x-y\|<2\varepsilon$. Thus, for every $x$ satisfying equation \eqref{eq2} we have 
		\begin{align*}
			\|Tx\|\le \|Ty\| + \|T\|\,\|x-y\|\leq \sup_{y\in M_G}\|Ty\|+2\|T\|\varepsilon,
		\end{align*}
		and consequently, $\|T\|_G \le \sup_{\|Gx\|>1-\delta}\|Tx\| \le \sup_{y\in M_G}\|Ty\|+2\|T\|\varepsilon$. Letting $\varepsilon\to0$ gives $\|T\|_G \le \sup_{x\in M_G}\|Tx\|$. Therefore, $\|T\|_G = \sup_{x\in M_G}\|Tx\|$. 
		
		
		\medskip
		
		
		\noindent
		$(iii)\implies (i).$ $(iii)\implies (i)$. Assume that $(i)$ fails. Then there exist $v_n\in S_H\cap E^\perp$ such that $\|Gv_n\|\to1$. Let $X=\overline{\operatorname{span}}\{v_n\}$ and denote by $P_X$ the orthogonal projection onto $X$. Since $P_Xv_n=v_n$, we have $\|P_Xv_n\|=1$, and because $\|Gv_n\|\to1$ it follows that $1\in S_G(P_X)$. Hence $\|P_X\|_G=1$. Since $X \subset E^\perp$, it follows that $E\subset X^\perp$. Because $M_G \subset E$, we obtain $M_G \subset X^\perp$. Hence $P_X x = 0$ for every $x \in M_G$. Therefore, $\sup_{x\in M_G}\|P_Xx\|=0,$ which contradicts $(iii)$. Thus $(i)$ must hold. 
		
		\medskip
		
		This completes the proof.
	\end{proof}

\end{document}
